\documentclass[11pt,aspectratio=169,xcolor=dvipsnames]{beamer}

\usetheme[
  outer/progressbar=foot,
]{metropolis}

\usepackage{amsmath}
\usepackage{amssymb}
\usepackage{mathtools}
\usepackage{booktabs}
\usepackage{xcolor}
\usepackage{tikz}
\usepackage{fontspec}
\usepackage{polyglossia}
\setmainlanguage{polish}

\IfFontExistsTF{Fira Sans Light}{
  \setsansfont[BoldFont={Fira Sans}]{Fira Sans Light}
}{
  \setsansfont{Arial}
}
\IfFontExistsTF{Fira Mono}{
  \setmonofont{Fira Mono}
}{
  \setmonofont{Menlo}
}

\setbeamercolor{palette primary}{bg=BlueViolet,fg=white}
\setbeamercolor{palette secondary}{bg=BlueViolet,fg=white}
\setbeamercolor{title separator}{bg=Red,fg=BlueViolet}
\setbeamercolor{progress bar}{bg=OrangeRed,fg=BlueViolet}
\beamertemplatenavigationsymbolsempty

\newcommand{\E}{\mathbb{E}}

\title{Modele VAR}
\subtitle{ASC 2026}
\author{Piotr Żoch}
\date{}

\begin{document}

\begin{frame}[plain]
  \vspace{1.2cm}
  {\usebeamerfont{title}\usebeamercolor[fg]{title}Modele VAR\par}
  \vspace{0.35cm}
  {\usebeamerfont{subtitle}\usebeamercolor[fg]{subtitle}ASC 2026\par}
  \vspace{0.9cm}
  {\usebeamerfont{author}Piotr Żoch\par}
\end{frame}

\begin{frame}{Plan}
  \begin{itemize}
    \item Model VAR
    \item Budowa modelu VAR
    \item Stabilność
    \item Estymacja i diagnostyka
    \item Prognozowanie, IRF i FEVD
  \end{itemize}
\end{frame}

\section{Model VAR}

\begin{frame}{Model VAR}
  \begin{itemize}
    \item VAR: \textbf{v}ector \textbf{a}uto\textbf{r}egression.
    \item Modeluje wspólną dynamikę wielu zmiennych bez narzucania z góry podziału na zmienne endogeniczne i egzogeniczne.
    \item Przydatny do:
      \begin{itemize}
        \item prognozowania szeregów czasowych,
        \item opisu zależności dynamicznych,
        \item analizy reakcji zmiennych na zaburzenia.
      \end{itemize}
    \item Przy dodatkowych założeniach identyfikacyjnych prowadzi do modeli SVAR.
  \end{itemize}
\end{frame}

\begin{frame}{Model VAR}
  Niech
  \[
    y_t = (y_{1,t}, y_{2,t}, \ldots, y_{n,t})'.
  \]
  Model VAR($p$):
  \[
    y_t = A_0 + A_1 y_{t-1} + A_2 y_{t-2} + \ldots + A_p y_{t-p} + \varepsilon_t,
  \]
  gdzie
  \[
    \varepsilon_t \sim (0,\Sigma).
  \]

  \begin{itemize}
    \item $A_0$ to wektor $n \times 1$.
    \item $A_1,\ldots,A_p$ oraz $\Sigma$ to macierze $n \times n$.
    \item Można dodać zmienne deterministyczne: stałą, trend, zmienne sezonowe, zmienne zero-jedynkowe.
  \end{itemize}
\end{frame}

\begin{frame}{Model VAR}
  Ten sam model można zapisać za pomocą wielomianu opóźnień:
  \[
    \left(I - A_1L - A_2L^2 - \ldots - A_pL^p\right)y_t = A_0 + \varepsilon_t,
  \]
  czyli
  \[
    A(L)y_t = A_0 + \varepsilon_t.
  \]

  \begin{itemize}
    \item Każde równanie ma te same zmienne objaśniające: opóźnienia wszystkich zmiennych z układu.
    \item Parametry pojedynczych równań zwykle nie są głównym obiektem interpretacji.
    \item Interesuje nas raczej cała dynamika systemu.
  \end{itemize}
\end{frame}

\begin{frame}{Dlaczego VAR?}
  \begin{itemize}
    \item Modele VAR powstały jako odpowiedź na problemy dużych modeli strukturalnych:
      \begin{itemize}
        \item arbitralny podział na zmienne endogeniczne i egzogeniczne,
        \item arbitralny wybór struktury dynamicznej,
        \item problemy z identyfikacją,
        \item słaba jakość prognoz w niektórych zastosowaniach.
      \end{itemize}
    \item Podejście VAR jest bardziej ateoretyczne:
      \begin{itemize}
        \item teoria pomaga wybrać zmienne,
        \item dane pomagają opisać zależności dynamiczne.
      \end{itemize}
  \end{itemize}
\end{frame}

\begin{frame}{Budowa modelu VAR}
  Dwie podstawowe decyzje:
  \begin{itemize}
    \item jakie zmienne uwzględnić,
    \item jaki rząd opóźnień $p$ wybrać.
  \end{itemize}

  \vspace{0.5cm}
  Dobór rzędu opóźnień:
  \begin{itemize}
    \item kryteria informacyjne: AIC, BIC, HQIC,
    \item testy istotności kolejnych opóźnień,
    \item diagnostyka reszt: brak autokorelacji i stabilna wariancja,
    \item sens ekonomiczny i częstotliwość danych.
  \end{itemize}
\end{frame}

\section{Stabilność}

\begin{frame}{VAR -- stacjonarność}
  Dla AR($p$) analizowaliśmy pierwiastki wielomianu opóźnień:
  \[
    A(L)y_t = a_0 + \varepsilon_t.
  \]

  Analogicznie w VAR($p$):
  \[
    A(L)y_t = A_0 + \varepsilon_t.
  \]

  \begin{itemize}
    \item VAR jest stabilny, jeśli pierwiastki równania
      \[
        \det A(z)=0
      \]
      leżą poza kołem jednostkowym.
    \item Intuicyjnie: wpływ jednorazowego szoku wygasa w czasie.
  \end{itemize}
\end{frame}

\begin{frame}{VAR(1) i wartości własne}
  Dla VAR(1):
  \[
    y_t = A_1 y_{t-1} + \varepsilon_t.
  \]

  Warunek stabilności:
  \[
    \det(I - A_1z)=0
  \]
  ma pierwiastki poza kołem jednostkowym.

  \begin{itemize}
    \item Równoważnie: wartości własne macierzy $A_1$ są w kole jednostkowym.
    \item Łatwo się pomylić:
      \begin{itemize}
        \item pierwiastki wielomianu opóźnień: poza kołem,
        \item wartości własne macierzy przejścia: w kole.
      \end{itemize}
  \end{itemize}
\end{frame}

\section{Estymacja i diagnostyka}

\begin{frame}{Estymacja i diagnostyka}
  \begin{itemize}
    \item Parametry VAR można estymować metodą najmniejszych kwadratów, równanie po równaniu.
    \item Po estymacji sprawdzamy:
      \begin{itemize}
        \item autokorelację reszt,
        \item stabilność modelu,
        \item istotność opóźnień,
        \item jakość prognoz poza próbą.
      \end{itemize}
    \item Autokorelacja reszt zwykle oznacza zbyt niski rząd opóźnień lub pominiętą strukturę deterministyczną.
  \end{itemize}
\end{frame}

\begin{frame}{Niestacjonarne zmienne w VAR}
  Co zrobić, jeśli zmienne w VAR są niestacjonarne?

  \begin{itemize}
    \item Oszacować VAR na poziomach:
      \begin{itemize}
        \item dobre dla prognoz i analizy zależności,
        \item ale klasyczna inferencja może być problematyczna.
      \end{itemize}
    \item Oszacować VAR na pierwszych różnicach:
      \begin{itemize}
        \item usuwa część niestacjonarności,
        \item ale traci informację o długookresowych relacjach.
      \end{itemize}
    \item Jeśli zmienne są skointegrowane: użyć VECM.
  \end{itemize}
\end{frame}

\section{Prognozy i impulsy}

\begin{frame}{Prognozowanie z VAR}
  \[
    y_t = A_0 + A_1y_{t-1} + \ldots + A_py_{t-p} + \varepsilon_t.
  \]

  Prognozy iteracyjne:
  \begin{align*}
    \E_T[y_{T+1}] &=
      A_0 + A_1y_T + A_2y_{T-1} + \ldots + A_py_{T-p+1},\\
    \E_T[y_{T+2}] &=
      A_0 + A_1\E_T[y_{T+1}] + A_2y_T + \ldots + A_py_{T-p+2}.
  \end{align*}

  \begin{itemize}
    \item W kolejnych krokach prognozy zastępują nieobserwowane przyszłe wartości.
    \item Przedziały prognoz wymagają wariancji błędów prognozy.
  \end{itemize}
\end{frame}

\begin{frame}{VAR i reprezentacja VMA}
  Dla stabilnego VAR istnieje reprezentacja VMA($\infty$):
  \[
    y_t = \mu + \sum_{i=0}^{\infty}\Theta_i\varepsilon_{t-i}.
  \]

  \begin{itemize}
    \item Macierze $\Theta_i$ opisują rozchodzenie się zaburzeń w czasie.
    \item Na tej reprezentacji opierają się:
      \begin{itemize}
        \item funkcje reakcji na impuls,
        \item dekompozycja wariancji błędów prognozy,
        \item analiza mechanizmów transmisji.
      \end{itemize}
    \item W modelu zredukowanym reszty mogą być współskorelowane, więc interpretacja szoków wymaga ostrożności.
  \end{itemize}
\end{frame}

\begin{frame}{Najważniejsze wnioski}
  \begin{itemize}
    \item VAR opisuje wspólną dynamikę wielu zmiennych.
    \item Dobór opóźnień i diagnostyka reszt są kluczowe.
    \item Stabilność oznacza wygasanie skutków szoków.
    \item Prognozy VAR otrzymujemy iteracyjnie.
    \item IRF i FEVD opierają się na reprezentacji VMA modelu VAR.
  \end{itemize}
\end{frame}

\end{document}
